% Options for packages loaded elsewhere
\PassOptionsToPackage{unicode}{hyperref}
\PassOptionsToPackage{hyphens}{url}
\documentclass[
]{article}
\usepackage{xcolor}
\usepackage{amsmath,amssymb}
\setcounter{secnumdepth}{-\maxdimen} % remove section numbering
\usepackage{iftex}
\ifPDFTeX
  \usepackage[T1]{fontenc}
  \usepackage[utf8]{inputenc}
  \usepackage{textcomp} % provide euro and other symbols
\else % if luatex or xetex
  \usepackage{unicode-math} % this also loads fontspec
  \defaultfontfeatures{Scale=MatchLowercase}
  \defaultfontfeatures[\rmfamily]{Ligatures=TeX,Scale=1}
\fi
\usepackage{lmodern}
\ifPDFTeX\else
  % xetex/luatex font selection
\fi
% Use upquote if available, for straight quotes in verbatim environments
\IfFileExists{upquote.sty}{\usepackage{upquote}}{}
\IfFileExists{microtype.sty}{% use microtype if available
  \usepackage[]{microtype}
  \UseMicrotypeSet[protrusion]{basicmath} % disable protrusion for tt fonts
}{}
\makeatletter
\@ifundefined{KOMAClassName}{% if non-KOMA class
  \IfFileExists{parskip.sty}{%
    \usepackage{parskip}
  }{% else
    \setlength{\parindent}{0pt}
    \setlength{\parskip}{6pt plus 2pt minus 1pt}}
}{% if KOMA class
  \KOMAoptions{parskip=half}}
\makeatother
\usepackage{color}
\usepackage{fancyvrb}
\newcommand{\VerbBar}{|}
\newcommand{\VERB}{\Verb[commandchars=\\\{\}]}
\DefineVerbatimEnvironment{Highlighting}{Verbatim}{commandchars=\\\{\}}
% Add ',fontsize=\small' for more characters per line
\newenvironment{Shaded}{}{}
\newcommand{\AlertTok}[1]{\textcolor[rgb]{1.00,0.00,0.00}{\textbf{#1}}}
\newcommand{\AnnotationTok}[1]{\textcolor[rgb]{0.38,0.63,0.69}{\textbf{\textit{#1}}}}
\newcommand{\AttributeTok}[1]{\textcolor[rgb]{0.49,0.56,0.16}{#1}}
\newcommand{\BaseNTok}[1]{\textcolor[rgb]{0.25,0.63,0.44}{#1}}
\newcommand{\BuiltInTok}[1]{\textcolor[rgb]{0.00,0.50,0.00}{#1}}
\newcommand{\CharTok}[1]{\textcolor[rgb]{0.25,0.44,0.63}{#1}}
\newcommand{\CommentTok}[1]{\textcolor[rgb]{0.38,0.63,0.69}{\textit{#1}}}
\newcommand{\CommentVarTok}[1]{\textcolor[rgb]{0.38,0.63,0.69}{\textbf{\textit{#1}}}}
\newcommand{\ConstantTok}[1]{\textcolor[rgb]{0.53,0.00,0.00}{#1}}
\newcommand{\ControlFlowTok}[1]{\textcolor[rgb]{0.00,0.44,0.13}{\textbf{#1}}}
\newcommand{\DataTypeTok}[1]{\textcolor[rgb]{0.56,0.13,0.00}{#1}}
\newcommand{\DecValTok}[1]{\textcolor[rgb]{0.25,0.63,0.44}{#1}}
\newcommand{\DocumentationTok}[1]{\textcolor[rgb]{0.73,0.13,0.13}{\textit{#1}}}
\newcommand{\ErrorTok}[1]{\textcolor[rgb]{1.00,0.00,0.00}{\textbf{#1}}}
\newcommand{\ExtensionTok}[1]{#1}
\newcommand{\FloatTok}[1]{\textcolor[rgb]{0.25,0.63,0.44}{#1}}
\newcommand{\FunctionTok}[1]{\textcolor[rgb]{0.02,0.16,0.49}{#1}}
\newcommand{\ImportTok}[1]{\textcolor[rgb]{0.00,0.50,0.00}{\textbf{#1}}}
\newcommand{\InformationTok}[1]{\textcolor[rgb]{0.38,0.63,0.69}{\textbf{\textit{#1}}}}
\newcommand{\KeywordTok}[1]{\textcolor[rgb]{0.00,0.44,0.13}{\textbf{#1}}}
\newcommand{\NormalTok}[1]{#1}
\newcommand{\OperatorTok}[1]{\textcolor[rgb]{0.40,0.40,0.40}{#1}}
\newcommand{\OtherTok}[1]{\textcolor[rgb]{0.00,0.44,0.13}{#1}}
\newcommand{\PreprocessorTok}[1]{\textcolor[rgb]{0.74,0.48,0.00}{#1}}
\newcommand{\RegionMarkerTok}[1]{#1}
\newcommand{\SpecialCharTok}[1]{\textcolor[rgb]{0.25,0.44,0.63}{#1}}
\newcommand{\SpecialStringTok}[1]{\textcolor[rgb]{0.73,0.40,0.53}{#1}}
\newcommand{\StringTok}[1]{\textcolor[rgb]{0.25,0.44,0.63}{#1}}
\newcommand{\VariableTok}[1]{\textcolor[rgb]{0.10,0.09,0.49}{#1}}
\newcommand{\VerbatimStringTok}[1]{\textcolor[rgb]{0.25,0.44,0.63}{#1}}
\newcommand{\WarningTok}[1]{\textcolor[rgb]{0.38,0.63,0.69}{\textbf{\textit{#1}}}}
\usepackage{longtable,booktabs,array}
\newcounter{none} % for unnumbered tables
\usepackage{calc} % for calculating minipage widths
% Correct order of tables after \paragraph or \subparagraph
\usepackage{etoolbox}
\makeatletter
\patchcmd\longtable{\par}{\if@noskipsec\mbox{}\fi\par}{}{}
\makeatother
% Allow footnotes in longtable head/foot
\IfFileExists{footnotehyper.sty}{\usepackage{footnotehyper}}{\usepackage{footnote}}
\makesavenoteenv{longtable}
\setlength{\emergencystretch}{3em} % prevent overfull lines
\providecommand{\tightlist}{%
  \setlength{\itemsep}{0pt}\setlength{\parskip}{0pt}}
\usepackage{bookmark}
\IfFileExists{xurl.sty}{\usepackage{xurl}}{} % add URL line breaks if available
\urlstyle{same}
\hypersetup{
  hidelinks,
  pdfcreator={LaTeX via pandoc}}

\author{}
\date{}

\begin{document}

\section{Lux Bridge Security Audit
Report}\label{lux-bridge-security-audit-report}

\textbf{Auditor}: Trail of Bits-Level Independent Review \textbf{Date}:
2025-01-30 \textbf{Scope}: \texttt{/contracts/bridge/} directory
\textbf{Severity Ratings}: Critical \textbar{} High \textbar{} Medium
\textbar{} Low \textbar{} Informational

\begin{center}\rule{0.5\linewidth}{0.5pt}\end{center}

\subsection{Executive Summary}\label{executive-summary}

This audit covers the Lux Bridge smart contract system, including
cross-chain token bridging, MPC oracle signature verification, vault
management, and yield strategy integration. The codebase contains
multiple bridge implementations with varying security postures.

\textbf{Overall Assessment}: The bridge system has several critical and
high-severity vulnerabilities that require immediate attention before
production deployment.

{\def\LTcaptype{none} % do not increment counter
\begin{longtable}[]{@{}ll@{}}
\toprule\noalign{}
Severity & Count \\
\midrule\noalign{}
\endhead
\bottomrule\noalign{}
\endlastfoot
Critical & 4 \\
High & 8 \\
Medium & 12 \\
Low & 9 \\
Informational & 6 \\
\end{longtable}
}

\begin{center}\rule{0.5\linewidth}{0.5pt}\end{center}

\subsection{Critical Findings}\label{critical-findings}

\subsubsection{C-01: XChainVault Burn Proof Verification is Effectively
Disabled}\label{c-01-xchainvault-burn-proof-verification-is-effectively-disabled}

\textbf{File}: \texttt{/contracts/bridge/XChainVault.sol} (Lines
304-312)

\textbf{Description}: The \texttt{\_verifyBurnProof} function contains a
critical vulnerability where burn proof verification is completely
bypassed. The function returns \texttt{true} if any proof data is
provided, regardless of content.

\begin{Shaded}
\begin{Highlighting}[]
\NormalTok{function \_verifyBurnProof(}
\NormalTok{    bytes32 vaultId,}
\NormalTok{    uint256 amount,}
\NormalTok{    bytes calldata proof}
\NormalTok{) private view returns (bool) \{}
\NormalTok{    // Verify Warp message signature}
\NormalTok{    // This would call the Warp precompile to verify the BLS signature}
\NormalTok{    // For now, simplified verification}
\NormalTok{    return proof.length \textgreater{} 0;  // CRITICAL: No actual verification!}
\NormalTok{\}}
\end{Highlighting}
\end{Shaded}

\textbf{Impact}: An attacker can drain all vaulted tokens by providing
any non-empty bytes as proof. Total loss of funds.

\textbf{Recommendation}: Implement proper Warp precompile verification:

\begin{Shaded}
\begin{Highlighting}[]
\NormalTok{function \_verifyBurnProof(...) private view returns (bool) \{}
\NormalTok{    (IWarp.WarpMessage memory message, bool valid) =}
\NormalTok{        IWarp(WARP\_PRECOMPILE).getVerifiedWarpMessage(0);}
\NormalTok{    if (!valid) return false;}
\NormalTok{    // Decode and verify message content matches expected burn}
\NormalTok{    return \_validateBurnMessage(message, vaultId, amount);}
\NormalTok{\}}
\end{Highlighting}
\end{Shaded}

\begin{center}\rule{0.5\linewidth}{0.5pt}\end{center}

\subsubsection{C-02: ecrecover Without Zero Address
Check}\label{c-02-ecrecover-without-zero-address-check}

\textbf{Files}:

\begin{itemize}
\tightlist
\item
  \texttt{/contracts/bridge/Bridge.sol} (Lines 319-328, legacy)
\item
  \texttt{/contracts/bridge/Teleport.sol} (Lines 271-273)
\end{itemize}

\textbf{Description}: The \texttt{recoverSigner} function uses raw
\texttt{ecrecover} without checking for the zero address return value.
When signature recovery fails, \texttt{ecrecover} returns
\texttt{address(0)} rather than reverting.

\begin{Shaded}
\begin{Highlighting}[]
\NormalTok{function recoverSigner(bytes32 message\_, bytes memory sig\_) internal pure returns (address) \{}
\NormalTok{    (uint8 v, bytes32 r, bytes32 s) = splitSignature(sig\_);}
\NormalTok{    return ecrecover(message\_, v, r, s);  // May return address(0)}
\NormalTok{\}}
\end{Highlighting}
\end{Shaded}

\textbf{Impact}: If \texttt{MPCOracleAddrMap{[}address(0){]}} is ever
set (intentionally or through upgrade mishap), arbitrary minting becomes
possible. Additionally, malformed signatures could cause unexpected
behavior.

\textbf{Recommendation}:

\begin{Shaded}
\begin{Highlighting}[]
\NormalTok{address signer = ecrecover(message\_, v, r, s);}
\NormalTok{require(signer != address(0), "Invalid signature");}
\NormalTok{return signer;}
\end{Highlighting}
\end{Shaded}

\begin{center}\rule{0.5\linewidth}{0.5pt}\end{center}

\subsubsection{C-03: Chain ID Not Validated in Teleporter Message
Signatures}\label{c-03-chain-id-not-validated-in-teleporter-message-signatures}

\textbf{File}: \texttt{/contracts/bridge/teleport/Teleporter.sol} (Lines
228-234)

\textbf{Description}: The \texttt{mintDeposit}
function\textquotesingle s message hash does not include the destination
chain ID (Lux chain), only the source chain ID. This allows cross-chain
replay attacks.

\begin{Shaded}
\begin{Highlighting}[]
\NormalTok{bytes32 messageHash = keccak256(abi.encodePacked(}
\NormalTok{    "DEPOSIT",}
\NormalTok{    srcChainId,        // Source chain}
\NormalTok{    depositNonce,}
\NormalTok{    recipient,}
\NormalTok{    amount}
\NormalTok{    // MISSING: block.chainid (destination chain)}
\NormalTok{));}
\end{Highlighting}
\end{Shaded}

\textbf{Impact}: A valid signature for minting on one Lux fork/testnet
can be replayed on another chain with the same contract address and
nonce state.

\textbf{Recommendation}: Include destination chain ID in message hash:

\begin{Shaded}
\begin{Highlighting}[]
\NormalTok{bytes32 messageHash = keccak256(abi.encodePacked(}
\NormalTok{    "DEPOSIT",}
\NormalTok{    srcChainId,}
\NormalTok{    block.chainid,     // Destination chain}
\NormalTok{    depositNonce,}
\NormalTok{    recipient,}
\NormalTok{    amount}
\NormalTok{));}
\end{Highlighting}
\end{Shaded}

\begin{center}\rule{0.5\linewidth}{0.5pt}\end{center}

\subsubsection{C-04: Withdraw Nonce Predictability Enables
Front-Running}\label{c-04-withdraw-nonce-predictability-enables-front-running}

\textbf{File}: \texttt{/contracts/bridge/teleport/Teleporter.sol} (Lines
318-325)

\textbf{Description}: The \texttt{burnForWithdraw} function generates
withdraw nonces using predictable on-chain values:

\begin{Shaded}
\begin{Highlighting}[]
\NormalTok{withdrawNonce = uint256(keccak256(abi.encodePacked(}
\NormalTok{    block.timestamp,}
\NormalTok{    msg.sender,}
\NormalTok{    amount,}
\NormalTok{    srcChainId,}
\NormalTok{    block.number}
\NormalTok{)));}
\end{Highlighting}
\end{Shaded}

\textbf{Impact}: An attacker can predict the nonce before transaction
confirmation, potentially enabling:

\begin{enumerate}
\def\labelenumi{\arabic{enumi}.}
\tightlist
\item
  Front-running to claim the withdrawal on the source chain before the
  legitimate user
\item
  Denial of service by pre-processing the nonce
\end{enumerate}

\textbf{Recommendation}: Use a monotonically increasing counter:

\begin{Shaded}
\begin{Highlighting}[]
\NormalTok{withdrawNonce = ++\_withdrawNonceCounter;}
\end{Highlighting}
\end{Shaded}

\begin{center}\rule{0.5\linewidth}{0.5pt}\end{center}

\subsection{High Severity Findings}\label{high-severity-findings}

\subsubsection{H-01: Missing Reentrancy Protection in Legacy Bridge
Contract}\label{h-01-missing-reentrancy-protection-in-legacy-bridge-contract}

\textbf{File}: \texttt{/contracts/bridge/Teleport.sol} (Lines 197-246)

\textbf{Description}: The \texttt{bridgeMintStealth} function lacks
reentrancy protection despite minting tokens (which could trigger
callbacks via hooks in some token implementations).

\begin{Shaded}
\begin{Highlighting}[]
\NormalTok{function bridgeMintStealth(...) public returns (address) \{  // No nonReentrant}
\NormalTok{    // ... validation ...}
\NormalTok{    varStruct.token.mint(payoutAddr, feeAmount);     // External call}
\NormalTok{    varStruct.token.mint(varStruct.toTargetAddr, netAmount);  // External call}
\NormalTok{    addMappingStealth(signedTXInfo);  // State update AFTER external calls}
\NormalTok{    // ...}
\NormalTok{\}}
\end{Highlighting}
\end{Shaded}

\textbf{Impact}: If the token implements ERC777 or has receive hooks, a
malicious actor could reenter and potentially double-mint.

\textbf{Recommendation}: Add \texttt{nonReentrant} modifier and follow
checks-effects-interactions pattern.

\begin{center}\rule{0.5\linewidth}{0.5pt}\end{center}

\subsubsection{H-02: TeleportProposalBridge Allows Message Execution
After
Expiry}\label{h-02-teleportproposalbridge-allows-message-execution-after-expiry}

\textbf{File}: \texttt{/contracts/bridge/TeleportProposalBridge.sol}
(Lines 405-422)

\textbf{Description}: The \texttt{executeMessage} function marks expired
messages as executed before reverting:

\begin{Shaded}
\begin{Highlighting}[]
\NormalTok{if (block.timestamp \textgreater{} message.timestamp + MESSAGE\_EXPIRY) \{}
\NormalTok{    message.executed = true; // Mark as executed (expired)}
\NormalTok{    emit MessageExpired(messageId);}
\NormalTok{    revert MessageExpiredError();}
\NormalTok{\}}
\end{Highlighting}
\end{Shaded}

\textbf{Impact}: While the transaction reverts, the state change
persists if this is called via a try/catch block, permanently blocking
legitimate re-execution attempts.

\textbf{Recommendation}: Do not modify state before reverting:

\begin{Shaded}
\begin{Highlighting}[]
\NormalTok{if (block.timestamp \textgreater{} message.timestamp + MESSAGE\_EXPIRY) \{}
\NormalTok{    revert MessageExpiredError();}
\NormalTok{\}}
\end{Highlighting}
\end{Shaded}

\begin{center}\rule{0.5\linewidth}{0.5pt}\end{center}

\subsubsection{H-03: Unbounded MPC Oracle
Set}\label{h-03-unbounded-mpc-oracle-set}

\textbf{Files}:

\begin{itemize}
\tightlist
\item
  \texttt{/contracts/bridge/Bridge.sol} (Lines 115-118)
\item
  \texttt{/contracts/bridge/teleport/Teleporter.sol} (Lines 404-416)
\end{itemize}

\textbf{Description}: MPC oracles can be added without bound and cannot
be fully revoked (mapping entries persist as \texttt{false} rather than
being deleted). There\textquotesingle s also no mechanism to enumerate
or audit current oracles.

\begin{Shaded}
\begin{Highlighting}[]
\NormalTok{function setMPCOracle(address MPCO\_) public onlyAdmin \{}
\NormalTok{    addMPCMapping(MPCO\_);  // Only sets to true, no way to enumerate or remove}
\NormalTok{\}}
\end{Highlighting}
\end{Shaded}

\textbf{Impact}:

\begin{enumerate}
\def\labelenumi{\arabic{enumi}.}
\tightlist
\item
  Stale or compromised oracle addresses may remain in mappings
\item
  No audit trail of oracle changes
\item
  Potential for validator collusion attacks if too many oracles are
  added
\end{enumerate}

\textbf{Recommendation}: Implement bounded oracle set with enumeration:

\begin{Shaded}
\begin{Highlighting}[]
\NormalTok{address[] public oracleList;}
\NormalTok{uint256 public constant MAX\_ORACLES = 10;}
\NormalTok{mapping(address =\textgreater{} uint256) public oracleIndex;}

\NormalTok{function addOracle(address oracle) external onlyAdmin \{}
\NormalTok{    require(oracleList.length \textless{} MAX\_ORACLES, "Max oracles");}
\NormalTok{    require(!mpcOracles[oracle], "Already oracle");}
\NormalTok{    oracleIndex[oracle] = oracleList.length;}
\NormalTok{    oracleList.push(oracle);}
\NormalTok{    mpcOracles[oracle] = true;}
\NormalTok{\}}
\end{Highlighting}
\end{Shaded}

\begin{center}\rule{0.5\linewidth}{0.5pt}\end{center}

\subsubsection{H-04: ETHVault Withdrawal Allowance
Bypass}\label{h-04-ethvault-withdrawal-allowance-bypass}

\textbf{File}: \texttt{/contracts/bridge/ETHVault.sol} (Lines 47-61)

\textbf{Description}: The withdrawal function has a typo in the error
message and potentially allows withdrawal without sufficient allowance
due to ordering:

\begin{Shaded}
\begin{Highlighting}[]
\NormalTok{function withdraw(uint256 amount\_, address receiver\_, address owner\_) external \{}
\NormalTok{    require(amount\_ \textless{}= address(this).balance, "Insufficient balance");}
\NormalTok{    if (msg.sender != owner\_) \{}
\NormalTok{        uint256 allowed = allowance(owner\_, msg.sender);}
\NormalTok{        require(allowed \textgreater{}= amount\_, "Invalid alowance");  // Typo, but functional}
\NormalTok{    \}}
\NormalTok{    \_burn(owner\_, amount\_);  // Burns shares}
\NormalTok{    (bool success, ) = payable(receiver\_).call\{value: amount\_\}("");}
\end{Highlighting}
\end{Shaded}

\textbf{Impact}: The allowance is checked but not decremented for
non-owner withdrawals, enabling unlimited withdrawals once any allowance
is granted.

\textbf{Recommendation}: Decrement allowance for non-owner:

\begin{Shaded}
\begin{Highlighting}[]
\NormalTok{if (msg.sender != owner\_) \{}
\NormalTok{    \_spendAllowance(owner\_, msg.sender, amount\_);}
\NormalTok{\}}
\end{Highlighting}
\end{Shaded}

\begin{center}\rule{0.5\linewidth}{0.5pt}\end{center}

\subsubsection{H-05: LiquidVault Strategy Manipulation via
Timestamp}\label{h-05-liquidvault-strategy-manipulation-via-timestamp}

\textbf{File}: \texttt{/contracts/bridge/teleport/LiquidVault.sol}
(Lines 187-214)

\textbf{Description}: Strategy allocation and deallocation signatures
include \texttt{block.timestamp}:

\begin{Shaded}
\begin{Highlighting}[]
\NormalTok{bytes32 messageHash = keccak256(abi.encodePacked(}
\NormalTok{    "ALLOCATE",}
\NormalTok{    strategyIndex,}
\NormalTok{    amount,}
\NormalTok{    block.timestamp  // Timestamp dependency}
\NormalTok{));}
\end{Highlighting}
\end{Shaded}

\textbf{Impact}: MPC oracle must sign messages with exact timestamp
matching. This creates a race condition where valid signatures become
invalid after the block containing them is mined. Miners can also
manipulate timestamps within bounds.

\textbf{Recommendation}: Use nonce-based signatures instead of
timestamps:

\begin{Shaded}
\begin{Highlighting}[]
\NormalTok{bytes32 messageHash = keccak256(abi.encodePacked(}
\NormalTok{    "ALLOCATE",}
\NormalTok{    strategyIndex,}
\NormalTok{    amount,}
\NormalTok{    allocationNonce++}
\NormalTok{));}
\end{Highlighting}
\end{Shaded}

\begin{center}\rule{0.5\linewidth}{0.5pt}\end{center}

\subsubsection{H-06: YieldBridgeVault Missing Access Control on
Harvest}\label{h-06-yieldbridgevault-missing-access-control-on-harvest}

\textbf{File}: \texttt{/contracts/bridge/yield/YieldBridgeVault.sol}
(Lines 216-236)

\textbf{Description}: The \texttt{harvestYield} function is publicly
callable without access control:

\begin{Shaded}
\begin{Highlighting}[]
\NormalTok{function harvestYield(address asset) external assetSupported(asset) returns (uint256 totalYield) \{}
\NormalTok{    // No access control modifier}
\end{Highlighting}
\end{Shaded}

\textbf{Impact}: Anyone can trigger harvests, potentially:

\begin{enumerate}
\def\labelenumi{\arabic{enumi}.}
\tightlist
\item
  Front-running yield distribution for MEV extraction
\item
  Griefing by triggering unnecessary gas expenditure
\item
  Timing manipulation for arbitrage
\end{enumerate}

\textbf{Recommendation}: Add access control or keeper pattern:

\begin{Shaded}
\begin{Highlighting}[]
\NormalTok{function harvestYield(address asset) external onlyKeeper assetSupported(asset) \{...\}}
\end{Highlighting}
\end{Shaded}

\begin{center}\rule{0.5\linewidth}{0.5pt}\end{center}

\subsubsection{H-07: Stale Backing Attestation
Bypass}\label{h-07-stale-backing-attestation-bypass}

\textbf{File}: \texttt{/contracts/bridge/teleport/Teleporter.sol} (Lines
484-499)

\textbf{Description}: The \texttt{\_checkBackingRatio} function silently
bypasses backing verification when attestation is stale:

\begin{Shaded}
\begin{Highlighting}[]
\NormalTok{if (block.timestamp {-} attestation.timestamp \textgreater{} 24 hours) \{}
\NormalTok{    // Stale attestation, allow minting but log warning}
\NormalTok{    // In production, might want stricter enforcement}
\NormalTok{    return;  // SILENTLY ALLOWS UNBACKED MINTING}
\NormalTok{\}}
\end{Highlighting}
\end{Shaded}

\textbf{Impact}: If MPC fails to update backing attestation for
\textgreater24 hours, unlimited unbacked tokens can be minted, breaking
the core invariant
\texttt{totalMinted\ \textless{}=\ totalBackingOnSourceChain}.

\textbf{Recommendation}: Revert on stale attestation in production:

\begin{Shaded}
\begin{Highlighting}[]
\NormalTok{if (block.timestamp {-} attestation.timestamp \textgreater{} 24 hours) \{}
\NormalTok{    revert StaleBackingAttestation();}
\NormalTok{\}}
\end{Highlighting}
\end{Shaded}

\begin{center}\rule{0.5\linewidth}{0.5pt}\end{center}

\subsubsection{H-08: LiquidETH Yield Index Integer
Overflow}\label{h-08-liquideth-yield-index-integer-overflow}

\textbf{File}: \texttt{/contracts/bridge/teleport/LiquidETH.sol} (Lines
303-314)

\textbf{Description}: The yield index calculation can overflow in edge
cases:

\begin{Shaded}
\begin{Highlighting}[]
\NormalTok{uint256 yieldPerDebt = amount * 1e18 / totalDebt;}
\NormalTok{yieldIndex += yieldPerDebt;  // Can overflow if many small yields accumulate}
\end{Highlighting}
\end{Shaded}

\textbf{Impact}: If \texttt{yieldIndex} overflows, all user debt
calculations become incorrect, potentially allowing debt escape or
incorrect liquidations.

\textbf{Recommendation}: Use checked math with upper bound:

\begin{Shaded}
\begin{Highlighting}[]
\NormalTok{require(yieldIndex + yieldPerDebt \textgreater{}= yieldIndex, "Yield index overflow");}
\NormalTok{require(yieldIndex + yieldPerDebt \textless{} type(uint128).max, "Yield index too large");}
\end{Highlighting}
\end{Shaded}

\begin{center}\rule{0.5\linewidth}{0.5pt}\end{center}

\subsection{Medium Severity Findings}\label{medium-severity-findings}

\subsubsection{M-01: Signature Malleability Not
Prevented}\label{m-01-signature-malleability-not-prevented}

\textbf{Files}: All contracts using raw ECDSA

\textbf{Description}: The signature splitting functions
don\textquotesingle t enforce low-S values, enabling signature
malleability (two valid signatures for same message).

\begin{Shaded}
\begin{Highlighting}[]
\NormalTok{function splitSignature(bytes memory sig) internal pure returns (uint8 v, bytes32 r, bytes32 s) \{}
\NormalTok{    // No check for s \textless{}= secp256k1n/2}
\NormalTok{\}}
\end{Highlighting}
\end{Shaded}

\textbf{Impact}: While replay protection uses the full signature as key
(mitigating direct replay), malleability can cause issues with off-chain
systems expecting unique signatures.

\textbf{Recommendation}: Use OpenZeppelin\textquotesingle s ECDSA
library which enforces low-S.

\begin{center}\rule{0.5\linewidth}{0.5pt}\end{center}

\subsubsection{M-02: Bridge.sol Fee Rate
Unbounded}\label{m-02-bridgesol-fee-rate-unbounded}

\textbf{File}: \texttt{/contracts/bridge/Bridge.sol} (Lines 26, 87-93)

\textbf{Description}: Fee rate can be set to any value without upper
bound:

\begin{Shaded}
\begin{Highlighting}[]
\NormalTok{uint256 public feeRate = 100; // 1\%}
\NormalTok{// ...}
\NormalTok{function setpayoutAddressess(address payoutAddress\_, uint256 feeRate\_) public onlyAdmin \{}
\NormalTok{    payoutAddress = payoutAddress\_;}
\NormalTok{    feeRate = feeRate\_;  // No upper bound check}
\NormalTok{\}}
\end{Highlighting}
\end{Shaded}

\textbf{Impact}: Admin can set 100\% fee, effectively stealing all
bridge transfers.

\textbf{Recommendation}: Add maximum fee cap:

\begin{Shaded}
\begin{Highlighting}[]
\NormalTok{require(feeRate\_ \textless{}= 1000, "Fee exceeds 10\%");  // Max 10\%}
\end{Highlighting}
\end{Shaded}

\begin{center}\rule{0.5\linewidth}{0.5pt}\end{center}

\subsubsection{M-03: Missing Zero Address
Validation}\label{m-03-missing-zero-address-validation}

\textbf{Files}: Multiple contracts

\textbf{Locations}:

\begin{itemize}
\tightlist
\item
  \texttt{Bridge.sol:setpayoutAddressess} - payoutAddress
\item
  \texttt{BridgeVault.sol:\_addNewERC20Vault} - asset validation only
  checks not zero for vault
\item
  \texttt{XChainVault.sol:releaseFromVault} - recipient
\end{itemize}

\textbf{Impact}: Setting critical addresses to zero can lock funds or
cause reverts.

\begin{center}\rule{0.5\linewidth}{0.5pt}\end{center}

\subsubsection{M-04: TeleportBridge.sol Missing Chain ID in Claim
ID}\label{m-04-teleportbridgesol-missing-chain-id-in-claim-id}

\textbf{File}: \texttt{/contracts/bridge/teleport/TeleportBridge.sol}
(Lines 213-224)

\textbf{Description}: The claim ID calculation doesn\textquotesingle t
include \texttt{block.chainid}:

\begin{Shaded}
\begin{Highlighting}[]
\NormalTok{claimId = keccak256(abi.encode(}
\NormalTok{    claim.burnTxHash,}
\NormalTok{    claim.logIndex,}
\NormalTok{    claim.token,}
\NormalTok{    claim.amount,}
\NormalTok{    claim.toChainId,}
\NormalTok{    // Missing: block.chainid}
\NormalTok{    // ...}
\NormalTok{));}
\end{Highlighting}
\end{Shaded}

\textbf{Impact}: Same claim ID on different chains with same bridge
deployment could cause confusion in monitoring/accounting systems.

\begin{center}\rule{0.5\linewidth}{0.5pt}\end{center}

\subsubsection{M-05: BridgeVault Approve Max Pattern
Risk}\label{m-05-bridgevault-approve-max-pattern-risk}

\textbf{File}: \texttt{/contracts/bridge/BridgeVault.sol} (Lines 85)

\textbf{Description}: Infinite approval to newly created vaults:

\begin{Shaded}
\begin{Highlighting}[]
\NormalTok{IERC20(asset\_).approve(newVaultAddress, type(uint256).max);}
\end{Highlighting}
\end{Shaded}

\textbf{Impact}: If vault contract is compromised, infinite tokens can
be drained.

\textbf{Recommendation}: Use incremental approvals or approve only
needed amounts.

\begin{center}\rule{0.5\linewidth}{0.5pt}\end{center}

\subsubsection{M-06: Teleporter Peg Check Uses Hardcoded
Value}\label{m-06-teleporter-peg-check-uses-hardcoded-value}

\textbf{File}: \texttt{/contracts/bridge/teleport/Teleporter.sol} (Lines
449-453)

\textbf{Description}: \texttt{getCurrentPeg()} returns hardcoded
BASIS\_POINTS:

\begin{Shaded}
\begin{Highlighting}[]
\NormalTok{function getCurrentPeg() public view returns (uint256) \{}
\NormalTok{    // In a real implementation, this would query a DEX oracle}
\NormalTok{    // For now, we assume 1:1 peg}
\NormalTok{    return BASIS\_POINTS;  // Always returns 10000}
\NormalTok{\}}
\end{Highlighting}
\end{Shaded}

\textbf{Impact}: Peg protection is ineffective. Bridge continues
operating during de-peg events.

\begin{center}\rule{0.5\linewidth}{0.5pt}\end{center}

\subsubsection{M-07: LiquidVault Strategy Array Not Bounded on
Iteration}\label{m-07-liquidvault-strategy-array-not-bounded-on-iteration}

\textbf{File}: \texttt{/contracts/bridge/teleport/LiquidVault.sol}
(Lines 262-267)

\textbf{Description}: Harvest iterates all strategies without gas limit
consideration:

\begin{Shaded}
\begin{Highlighting}[]
\NormalTok{for (uint256 i = 0; i \textless{} strategies.length; i++) \{}
\NormalTok{    if (strategies[i].active) \{}
\NormalTok{        IYieldStrategy(strategies[i].adapter).harvest();}
\end{Highlighting}
\end{Shaded}

\textbf{Impact}: With MAX\_STRATEGIES=10, if each harvest is
gas-intensive, function could exceed block gas limit.

\begin{center}\rule{0.5\linewidth}{0.5pt}\end{center}

\subsubsection{M-08: ProposalBridge pendingMessageIds Array Grows
Unbounded}\label{m-08-proposalbridge-pendingmessageids-array-grows-unbounded}

\textbf{File}: \texttt{/contracts/bridge/TeleportProposalBridge.sol}
(Lines 118, 378-380)

\textbf{Description}: Messages are added to \texttt{pendingMessageIds}
but never removed:

\begin{Shaded}
\begin{Highlighting}[]
\NormalTok{bytes32[] public pendingMessageIds;}
\NormalTok{// ...}
\NormalTok{pendingMessageIds.push(messageId);  // Only grows}
\end{Highlighting}
\end{Shaded}

\textbf{Impact}: Over time, view functions iterating this array will
exceed gas limits.

\begin{center}\rule{0.5\linewidth}{0.5pt}\end{center}

\subsubsection{M-09: Missing Events for Critical State
Changes}\label{m-09-missing-events-for-critical-state-changes}

\textbf{Files}: Multiple

\textbf{Locations}:

\begin{itemize}
\tightlist
\item
  \texttt{Bridge.sol:setWithdrawalEnabled} - no event
\item
  \texttt{ETHVault.sol} - no event for allowance changes
\item
  \texttt{LiquidVault.sol:setBufferBps} - has event but could be
  front-run
\end{itemize}

\begin{center}\rule{0.5\linewidth}{0.5pt}\end{center}

\subsubsection{M-10: Teleporter mintYield Missing Backing
Check}\label{m-10-teleporter-mintyield-missing-backing-check}

\textbf{File}: \texttt{/contracts/bridge/teleport/Teleporter.sol} (Lines
260-293)

\textbf{Description}: \texttt{mintYield} doesn\textquotesingle t call
\texttt{\_checkBackingRatio}:

\begin{Shaded}
\begin{Highlighting}[]
\NormalTok{function mintYield(...) external nonReentrant whenNotPaused \{}
\NormalTok{    // No checkPeg modifier}
\NormalTok{    // No \_checkBackingRatio call}
\end{Highlighting}
\end{Shaded}

\textbf{Impact}: Yield minting could violate backing invariant during
stress conditions.

\begin{center}\rule{0.5\linewidth}{0.5pt}\end{center}

\subsubsection{M-11: LRC20B Burns From Arbitrary
Accounts}\label{m-11-lrc20b-burns-from-arbitrary-accounts}

\textbf{File}: \texttt{/contracts/bridge/LRC20B.sol} (Lines 82-89)

\textbf{Description}: Admin can burn from any account without approval:

\begin{Shaded}
\begin{Highlighting}[]
\NormalTok{function bridgeBurn(address account, uint256 amount) public onlyAdmin returns (bool) \{}
\NormalTok{    \_burn(account, amount);  // No approval check}
\end{Highlighting}
\end{Shaded}

\textbf{Impact}: While intended for bridge operations, this is a
centralization risk. Admin can burn user tokens arbitrarily.

\begin{center}\rule{0.5\linewidth}{0.5pt}\end{center}

\subsubsection{M-12: YieldBridgeVault Strategy Weight
Validation}\label{m-12-yieldbridgevault-strategy-weight-validation}

\textbf{File}: \texttt{/contracts/bridge/yield/YieldBridgeVault.sol}
(Lines 327-348)

\textbf{Description}: Strategy weights aren\textquotesingle t validated
to sum to 100\%:

\begin{Shaded}
\begin{Highlighting}[]
\NormalTok{require(targetWeight \textless{}= BASIS\_POINTS, "YieldBridgeVault: invalid weight");}
\NormalTok{// No check that sum of all weights == BASIS\_POINTS}
\end{Highlighting}
\end{Shaded}

\textbf{Impact}: Under-allocated funds remain idle, over-allocated
causes revert on rebalance.

\begin{center}\rule{0.5\linewidth}{0.5pt}\end{center}

\subsection{Low Severity Findings}\label{low-severity-findings}

\subsubsection{L-01: Floating Pragma}\label{l-01-floating-pragma}

All contracts use \texttt{\^{}0.8.31} or similar. Pin to exact version
for reproducible builds.

\subsubsection{L-02: Missing NatSpec
Documentation}\label{l-02-missing-natspec-documentation}

Many functions lack @param and @return documentation.

\subsubsection{L-03: Unused Imports}\label{l-03-unused-imports}

\texttt{Bridge.sol} imports Strings twice (line 14, 17).

\subsubsection{L-04: Magic Numbers}\label{l-04-magic-numbers}

Fee calculations use raw numbers like \texttt{10\ **\ 4},
\texttt{10\ **\ 18}. Define named constants.

\subsubsection{L-05: Inconsistent
Naming}\label{l-05-inconsistent-naming}

\begin{itemize}
\tightlist
\item
  \texttt{setpayoutAddressess} (typo, should be
  \texttt{setPayoutAddresses})
\item
  Mixed camelCase and snake\_case in parameters
\end{itemize}

\subsubsection{L-06: Shadow Variables}\label{l-06-shadow-variables}

\texttt{TeleportVault.withdrawNonce} shadows state with parameter
\texttt{\_withdrawNonce}.

\subsubsection{L-07: State Variables
Visibility}\label{l-07-state-variables-visibility}

\texttt{withdrawalEnabled} in Bridge.sol lacks explicit visibility
(defaults to internal).

\subsubsection{L-08: Missing Return Value
Check}\label{l-08-missing-return-value-check}

ERC20 transfers in legacy Bridge.sol don\textquotesingle t use
SafeERC20:

\begin{Shaded}
\begin{Highlighting}[]
\NormalTok{IERC20(tokenAddr\_).transferFrom(msg.sender, address(vault), amount\_);}
\end{Highlighting}
\end{Shaded}

\subsubsection{L-09: Gas Inefficient Storage
Reads}\label{l-09-gas-inefficient-storage-reads}

Multiple storage reads in loops without caching in memory.

\begin{center}\rule{0.5\linewidth}{0.5pt}\end{center}

\subsection{Informational Findings}\label{informational-findings}

\subsubsection{I-01: Consider Using EIP-712 Typed Signatures
Everywhere}\label{i-01-consider-using-eip-712-typed-signatures-everywhere}

TeleportBridge.sol uses EIP-712 properly. Extend this to all signature
verification for better UX and replay protection.

\subsubsection{I-02: Centralization
Risks}\label{i-02-centralization-risks}

\begin{itemize}
\tightlist
\item
  Single admin can pause all operations
\item
  MPC oracle set controlled by admin
\item
  Fee rates controlled by admin
\end{itemize}

Consider timelock and multi-sig for admin functions.

\subsubsection{I-03: Lack of Emergency
Functions}\label{i-03-lack-of-emergency-functions}

No circuit breaker for individual functions. Consider granular pause
controls.

\subsubsection{I-04: Missing Getter
Functions}\label{i-04-missing-getter-functions}

\begin{itemize}
\tightlist
\item
  No way to enumerate all MPC oracles
\item
  No way to get all pending messages efficiently
\end{itemize}

\subsubsection{I-05: Upgrade Path Not
Defined}\label{i-05-upgrade-path-not-defined}

Contracts are not upgradeable. Plan migration path for critical fixes.

\subsubsection{I-06: Test Coverage
Unknown}\label{i-06-test-coverage-unknown}

No test files visible in scope. Recommend \textgreater95\% coverage for
bridge contracts.

\begin{center}\rule{0.5\linewidth}{0.5pt}\end{center}

\subsection{Validator Collusion
Analysis}\label{validator-collusion-analysis}

The MPC-based bridge architecture relies on threshold signatures from
validators. Analysis:

\subsubsection{Attack Vectors}\label{attack-vectors}

\begin{enumerate}
\def\labelenumi{\arabic{enumi}.}
\item
  \textbf{Colluding Majority}: If threshold validators collude, they can
  mint arbitrary tokens or sign fraudulent messages.
\item
  \textbf{Key Compromise}: Single validator key compromise allows
  signing share, combined with social engineering of others.
\item
  \textbf{Slashing Absent}: No economic penalty for malicious signing.
\end{enumerate}

\subsubsection{Mitigations Observed}\label{mitigations-observed}

\begin{itemize}
\tightlist
\item
  Threshold requirement in TeleportProposalBridge
\item
  Role-based access control
\item
  Nonce-based replay protection
\end{itemize}

\subsubsection{Recommendations}\label{recommendations}

\begin{enumerate}
\def\labelenumi{\arabic{enumi}.}
\tightlist
\item
  Implement validator bond/slash mechanism
\item
  Add time delays for large value transfers
\item
  Implement fraud proof system
\item
  Consider optimistic rollup pattern for additional security layer
\end{enumerate}

\begin{center}\rule{0.5\linewidth}{0.5pt}\end{center}

\subsection{Double-Spend Scenario
Analysis}\label{double-spend-scenario-analysis}

\subsubsection{Identified Vectors}\label{identified-vectors}

\begin{enumerate}
\def\labelenumi{\arabic{enumi}.}
\item
  \textbf{Cross-Chain Race Condition}: Burn on Lux + claim on source
  chain. If source chain release is faster than burn finality,
  double-spend possible.
\item
  \textbf{Reorg Vulnerability}: Deep reorg on source chain after Lux
  mint could result in double-spend.
\item
  \textbf{Nonce Prediction}: Predictable withdraw nonces (Teleporter)
  enable front-running.
\end{enumerate}

\subsubsection{Mitigations Required}\label{mitigations-required}

\begin{enumerate}
\def\labelenumi{\arabic{enumi}.}
\tightlist
\item
  Wait for sufficient block confirmations before minting
\item
  Implement challenge period for withdrawals
\item
  Use monotonic counters instead of hash-based nonces
\end{enumerate}

\begin{center}\rule{0.5\linewidth}{0.5pt}\end{center}

\subsection{Recommendations Summary}\label{recommendations-summary}

\subsubsection{Immediate Actions (Before
Mainnet)}\label{immediate-actions-before-mainnet}

\begin{enumerate}
\def\labelenumi{\arabic{enumi}.}
\tightlist
\item
  Fix XChainVault burn proof verification (C-01)
\item
  Add ecrecover zero address check (C-02)
\item
  Include destination chain ID in all signatures (C-03)
\item
  Replace predictable nonces with counters (C-04)
\item
  Add reentrancy guards to all state-modifying functions (H-01)
\item
  Fix ETHVault allowance handling (H-04)
\item
  Remove stale attestation bypass (H-07)
\end{enumerate}

\subsubsection{Short-Term (Next Release)}\label{short-term-next-release}

\begin{enumerate}
\def\labelenumi{\arabic{enumi}.}
\tightlist
\item
  Implement bounded oracle management (H-03)
\item
  Add fee rate caps (M-02)
\item
  Implement proper peg oracle (M-06)
\item
  Add comprehensive events (M-09)
\end{enumerate}

\subsubsection{Long-Term}\label{long-term}

\begin{enumerate}
\def\labelenumi{\arabic{enumi}.}
\tightlist
\item
  Add economic security (validator staking/slashing)
\item
  Implement fraud proofs
\item
  Consider upgradeability pattern
\item
  Add formal verification for core invariants
\end{enumerate}

\begin{center}\rule{0.5\linewidth}{0.5pt}\end{center}

\subsection{Files Reviewed}\label{files-reviewed}

{\def\LTcaptype{none} % do not increment counter
\begin{longtable}[]{@{}llllll@{}}
\toprule\noalign{}
File & Lines & Critical & High & Medium & Low \\
\midrule\noalign{}
\endhead
\bottomrule\noalign{}
\endlastfoot
Bridge.sol & 582 & 1 & 1 & 2 & 3 \\
Teleport.sol & 284 & 1 & 1 & 1 & 1 \\
Teleporter.sol & 521 & 2 & 2 & 2 & 1 \\
TeleportBridge.sol & 384 & 0 & 0 & 1 & 0 \\
BridgeVault.sol & 221 & 0 & 0 & 1 & 0 \\
XChainVault.sol & 358 & 1 & 0 & 0 & 0 \\
TeleportVault.sol & 300 & 0 & 0 & 0 & 1 \\
LiquidVault.sol & 432 & 0 & 2 & 1 & 1 \\
ETHVault.sol & 62 & 0 & 1 & 0 & 0 \\
LRC20B.sol & 97 & 0 & 0 & 1 & 0 \\
LiquidETH.sol & 571 & 0 & 1 & 0 & 1 \\
YieldBridgeVault.sol & 556 & 0 & 1 & 1 & 0 \\
TeleportProposalBridge.sol & 586 & 0 & 1 & 2 & 1 \\
\end{longtable}
}

\begin{center}\rule{0.5\linewidth}{0.5pt}\end{center}

\textbf{Audit Completed}: 2025-01-30 \textbf{Next Review Recommended}:
After all critical/high issues addressed

\end{document}
